\chapter{Conclusion}
\label{chap:conclusion}

Rust has become a very important consideration in the systems programming field. It has grown in demand over the last few years and has yet to show signs of slowing down. It still holds its title of the most admired programming languages as suggested by Stackoverflow\cite[]{Technology2024Stack}, only beaten by C++11 back in 2015. With ever-growing demand for memory safety in systems programming, Rust stands as a likely future competitor to the industry standard C++. Many industry giants are adopting Rust into their technology portfolio as described in chapter~\ref{chap:real_world} on real world adoption. 

However, Rust still has far to go in terms of its adoption. It is often criticized~\ref{chap:development} for its steep learning curve and restrictiveness when it comes to the borrow checker. Rust forces the developer to reformulate the way they think about ownership of data, compared to long-standing approaches favored by languages like C and C++. However, it is much more approachable by developers, as it is closer to C++ than functional programming languages like Haskell. Furthermore, most respondents would argue that they would be willing to use Rust, if the industry and the job market would catch up.

\paragraph{}

Comparing Rust and C++ performance yields very similar metrics when it comes to raw performance, Rust occasionally outperforming C++ code, and Rust even showing advantages when it comes to memory usage, even in comparison to C. However, code performance always depends on the implementation of said code, making head-to-head comparison of Rust and C++ difficult. More research should be done on comparing rewrites of systems from C++ to Rust.